% Owner: Iris (web app, live mode, mini-games, study logging).
\section{Application and User Study Design}
\label{sec:app}

ChromaFix is delivered as a web application so that a user can be screened, see how their
vision is modelled, and use the recoloring on their own images and camera without installing
anything. A FastAPI server wraps the \texttt{chromafix} Python package, and a React and
TypeScript front end talks to it through a single typed API client. All color math runs in
Python through the shared conversion functions of \cref{sec:simulation}, with one exception, the live
camera shader described below. The server processes uploaded images in memory and never
writes them to disk.

\paragraph{Screening and profile.}
The screening page shows 24 plates from \cref{sec:simulation}: three control plates that every
observer should read, and the rest spread evenly over the three deficiency types and the
difficulty levels. For each plate the user types the digit they see or selects ``I don't see a
number''. The answers are scored on the server into a profile $(t, s, c)$ (type, severity,
confidence). That profile is kept in the app's state and sent with every later request, so the
image, live and game pages all adapt to the same viewer. A control bar lets the user, or a
demonstrator, override the type and severity by hand.

\paragraph{Image mode.}
The user uploads an image and sees four panels side by side: the original, the original as
simulated for their profile, the recolored image, and the recolored image as simulated for
their profile, \ie $S\,R\,\mathbf{x}$ with recoloring applied first. The last panel is the one
that matters for a viewer with CVD; the second and fourth together let a viewer with normal
vision see what the recoloring changes. The recoloring method is chosen from the methods the
server reports as available, so methods that are not ready (the neural model, or personalized
parameters that have not been tuned yet) are hidden rather than failing.

\paragraph{Live mode.}
Live mode applies a method to webcam video in real time. Because simulation and the
matrix-based recoloring methods are each a single $3{\times}3$ matrix on linear RGB, a frame
can be processed entirely on the GPU. A WebGL~1 fragment shader decodes each sRGB pixel to
linear RGB with the IEC~61966-2-1 curve, multiplies it by the matrix, clamps the result to
$[0,1]$, and encodes it back to sRGB. The matrices come from the server, so the live view uses
exactly the same values as the Python pipeline. The server returns them row-major while GLSL
expects column-major and WebGL~1 does not allow the transpose flag, so the client transposes
before upload. For the ``recolored, as you see it'' view the client composes $S\,R$ into one
matrix. This clamps once instead of after each step, so it can differ slightly from the image
pipeline for pixels that recoloring pushes out of gamut. The page shows the achieved frame rate;
on our laptop it held 120~fps in every view, so the per-pixel matrix adds no visible cost.
Neural recoloring has no matrix form and is therefore not offered in live mode.

\paragraph{Color-match game.}
To measure whether recoloring helps people in practice, and not only in CIEDE2000, we built a
short game on the annotated test images (\cref{sec:data}). Each of its 12 rounds shows one
image with the two regions of an annotated pair outlined and labelled A and B, together with a
question that only color can answer, such as ``which wire is red?'' or ``which bananas are
ripe?''. The player clicks A or B as quickly as they can. The rounds cover charts, color-coded
items, food, maps and traffic signals, mostly red--green distinctions. The correct answer is A
in six rounds and B in six. Only test-split images are used, because the personalized method
was tuned on the tune split; the server rejects any game entry that points to a tune image or
to an unannotated pair. The image is recolored with the round's condition and, by default,
also simulated for the player's profile. With the profile set to a deficiency, players with
normal vision then perform the task under simulated CVD, while players with CVD can switch the
simulation off and use their own vision. Response time is measured from the moment the image
has loaded until the click.

\paragraph{Counterbalancing.}
Let $C$ be the list of conditions and $p$ the participant number. Round $k$ of the answer key is
shown under condition $C[(k + p) \bmod |C|]$, so each image rotates through all conditions
across consecutive participants, and every participant sees every condition about equally
often. The order of the rounds is then shuffled with a seed derived from the project seed and
$p$. The same participant therefore always receives the same sequence, which makes sessions
reproducible, while different participants see the rounds in different orders. $C$ contains
every matrix-based method the server can run; at the time of writing these are \emph{none}
and \emph{fixed}. The personalized method is added automatically once its
parameters are available. The neural method is excluded because it is slow on a laptop CPU.

\paragraph{Logging and privacy.}
Each answer is posted to the server and appended as one JSON line to a per-session log. A line
holds the session and participant IDs, the profile used, the game, round, image and condition,
the answer, whether it was correct, the response time in milliseconds, and a timestamp. These
logs are the input to the analysis in \cref{sec:study}. Participant IDs are anonymous codes of
the form \texttt{P07}. The server rejects any log line with fields outside this schema, so
names, e-mail addresses or photographs cannot enter the logs by accident, and it stores no
uploaded images.

\paragraph{Sessions.}
Because of the project timeline, the game is run in short demonstration sessions rather than
a formally recruited study. Each session runs on one laptop in a desktop browser. Before
playing, the participant is told that the session is anonymous, that only their answers and
response times are recorded under a code instead of their name, and that they may stop at any
time. A session consists of an optional screening, the 12 game rounds, and the per-condition
summary the game shows at the end. For participants with normal vision the profile is set to
deutan at severity 1.0, the most common deficiency~\cite{birch2012}, which matches the mostly
red--green rounds; participants whose screening finds a CVD keep their own profile.
% Participant counts come from the real logs via results/report/study_numbers.tex
% (scripts/analyze_study.py); never type them by hand.
\ifnum\StudyParticipants>0
\StudyParticipants{} people took part in \StudySessions{} sessions. Because the log records
the profile that was applied rather than the screening result, players with a simulated
deficiency cannot be told apart from players with a real one in the logs.
\else
No sessions had been logged at the time of writing.
\fi
With this few participants the results in \cref{sec:study} are indicative only, and a
recruited study that includes more participants with CVD is left to future work.
